\documentclass[11pt,a4paper]{article}
\usepackage{amsmath}
\usepackage{amssymb}
\usepackage{graphicx}
\usepackage{booktabs}
\usepackage{float}
\usepackage{geometry}
\usepackage{fancyhdr}

\geometry{margin=1in}
\pagestyle{fancy}
\fancyhf{}
\rhead{Supplementary Materials}
\lhead{VanaVaid}
\cfoot{\thepage}

\title{\textbf{VanaVaid: Supplementary Technical Materials}}
\author{AI Research Team}
\date{\today}

\begin{document}

\maketitle

\section{S1. Detailed Plant Compound Database}

\subsection{S1.1 Phytochemical Baseline Concentrations}

The following table documents literature-calibrated baseline phytochemical concentrations (\% dry weight) for all 15 medicinal plants:

\begin{table}[H]
\centering
\caption{Baseline Phytochemical Concentrations for 15 Indian Medicinal Plants}
\label{tab:phyto_baseline}
\small
\begin{tabular}{|l|l|l|l|l|l|}
\hline
\textbf{Plant Name} & \textbf{Alkaloid \%} & \textbf{Flavonoid \%} & \textbf{Terpenoid \%} & \textbf{Glycoside \%} & \textbf{Plant Family} \\
\hline
Turmeric & 3.5 ± 0.6 & 5.2 ± 0.8 & 2.3 ± 0.4 & 1.2 ± 0.3 & Zingiberaceae \\
\hline
Ashwagandha & 2.0 ± 0.5 & 3.8 ± 0.6 & 1.1 ± 0.3 & 2.5 ± 0.7 & Solanaceae \\
\hline
Neem & 1.5 ± 0.4 & 4.5 ± 0.7 & 2.8 ± 0.5 & 1.8 ± 0.4 & Meliaceae \\
\hline
Tulsi & 1.3 ± 0.3 & 6.2 ± 0.9 & 2.0 ± 0.4 & 1.0 ± 0.2 & Lamiaceae \\
\hline
Ginger & 1.0 ± 0.3 & 4.0 ± 0.6 & 3.0 ± 0.5 & 0.7 ± 0.2 & Zingiberaceae \\
\hline
Brahmi & 2.3 ± 0.5 & 4.8 ± 0.7 & 0.8 ± 0.2 & 3.5 ± 0.8 & Plantaginaceae \\
\hline
Amla & 0.6 ± 0.2 & 9.0 ± 1.2 & 0.4 ± 0.1 & 2.2 ± 0.5 & Phyllanthaceae \\
\hline
Garlic & 0.8 ± 0.2 & 2.8 ± 0.5 & 1.6 ± 0.3 & 1.5 ± 0.4 & Amaryllidaceae \\
\hline
Licorice & 0.5 ± 0.1 & 3.2 ± 0.5 & 1.4 ± 0.3 & 15.0 ± 2.0 & Fabaceae \\
\hline
Triphala & 1.2 ± 0.3 & 7.5 ± 1.0 & 0.9 ± 0.2 & 5.3 ± 0.9 & Mixed \\
\hline
\end{tabular}
\end{table}

\subsection{S1.2 Geographic Regional Characteristics}

\begin{table}[H]
\centering
\caption{Detailed Geographic and Agroclimatic Characteristics}
\label{tab:regions_detailed}
\small
\begin{tabular}{|l|l|l|l|l|l|}
\hline
\textbf{Region} & \textbf{Latitude} & \textbf{Temp (°C)} & \textbf{Rainfall (mm)} & \textbf{Soil pH} & \textbf{Altitude (m)} \\
\hline
Nilgiris, TN & 11.4°N & 20-25 & 1,800-2,000 & 6.0-6.8 & 2,600 \\
\hline
Himachal Pradesh & 31.7°N & 5-20 & 1,200-2,500 & 6.5-7.5 & 1,000-2,200 \\
\hline
Western Ghats & 13°N & 15-28 & 2,000-3,000 & 5.5-6.5 & 600-2,000 \\
\hline
Deccan Plateau & 17°N & 22-32 & 600-1,000 & 7.0-8.0 & 400-800 \\
\hline
Malwa Region & 23°N & 20-35 & 800-1,200 & 7.5-8.5 & 400-600 \\
\hline
Rajasthan & 28°N & 20-40 & 200-400 & 8.0-9.0 & 200-400 \\
\hline
Punjab & 31°N & 5-35 & 400-800 & 7.0-8.5 & 200-300 \\
\hline
Uttar Pradesh & 27°N & 5-40 & 600-1,200 & 7.5-8.5 & 60-400 \\
\hline
\end{tabular}
\end{table}

\section{S2. Hyperparameter Configuration}

\subsection{S2.1 Module 1 Gaussian Process Parameters}

\begin{itemize}
    \item \textbf{Kernel:} RBF + White Noise
    \item \textbf{RBF Signal Variance:} $\sigma^2 = 1.2$
    \item \textbf{RBF Length Scale:} $\ell = [0.15, 0.18, 0.22, 0.10, 0.08, 0.12]$ (per dimension)
    \item \textbf{Noise Variance:} $\sigma_n^2 = 0.05$
    \item \textbf{Optimizer:} L-BFGS-B
    \item \textbf{Max iterations:} 1,000
    \item \textbf{Learning rate (for gradient):} 0.01
\end{itemize}

\subsection{S2.2 Module 2 Bayesian Neural Network Parameters}

\begin{itemize}
    \item \textbf{Architecture:} Dense(128) → ReLU → Dropout(0.3) → Dense(64) → ReLU → Dropout(0.3) → Dense(32) → ReLU → Dropout(0.3) → Dense(1)
    \item \textbf{Optimizer:} Adam with $\beta_1=0.9$, $\beta_2=0.999$, $\epsilon=1e-7$
    \item \textbf{Learning rate:} 0.001
    \item \textbf{Batch size:} 16
    \item \textbf{Epochs:} 50
    \item \textbf{Dropout rate (MC):} 0.3
    \item \textbf{MC forward passes:} 30
    \item \textbf{L2 penalty:} 0.0001
\end{itemize}

\subsection{S2.3 Module 3 Graph Neural Network Parameters}

\begin{itemize}
    \item \textbf{Node embedding method:} DeepWalk
    \item \textbf{Embedding dimension:} 64
    \item \textbf{Walk length:} 40
    \item \textbf{Walks per node:} 10
    \item \textbf{Window size:} 5
    \item \textbf{Link predictor MLP:} Dense(256) → ReLU → Dense(128) → ReLU → Dense(64) → ReLU → Dense(1) → Sigmoid
    \item \textbf{Loss:} Binary cross-entropy
    \item \textbf{Optimizer:} Adam ($\eta = 0.001$)
    \item \textbf{Batch size:} 32
    \item \textbf{Epochs:} 100
\end{itemize}

\subsection{S2.4 Module 5 CNN Architecture Details}

\begin{verbatim}
Input: (224, 224, 5) multichannel image (RGB + NIR + UV)

Layer 1: Conv2D(32, kernel_size=3, padding='same')
         BatchNormalization
         ReLU
         MaxPooling2D(pool_size=2)
         → Output: (112, 112, 32)

Layer 2: Conv2D(64, kernel_size=3, padding='same')
         BatchNormalization
         ReLU
         MaxPooling2D(pool_size=2)
         → Output: (56, 56, 64)

Layer 3: Conv2D(128, kernel_size=3, padding='same')
         BatchNormalization
         ReLU
         MaxPooling2D(pool_size=2)
         → Output: (28, 28, 128)

Layer 4: GlobalAveragePooling2D()
         → Output: (128,)

Layer 5: Dense(128)
         ReLU
         Dropout(0.2)

Layer 6: Dense(4)
         Sigmoid (for normalized concentration outputs)
         → Output: [c_alk, c_flav, c_terp, c_glyc]
\end{verbatim}

\section{S3. Training Data Statistics}

\subsection{S3.1 Dataset Composition}

\begin{table}[H]
\centering
\caption{Dataset Split and Composition}
\label{tab:data_stats}
\begin{tabular}{|l|l|l|}
\hline
\textbf{Metric} & \textbf{Phytochemical Module (1)} & \textbf{Spectral Module (5)} \\
\hline
Total samples & 1,474 & 2,850 \\
\hline
Training set & 1,179 (80\%) & 1,995 (70\%) \\
\hline
Validation set & 148 (10\%) & 428 (15\%) \\
\hline
Test set & 147 (10\%) & 427 (15\%) \\
\hline
\end{tabular}
\end{table}

\subsection{S3.2 Feature Normalization}

All features normalized to zero mean and unit variance:
\[
\mathbf{x}_{\text{norm}} = \frac{\mathbf{x} - \mu}{\sigma}
\]

where $\mu$ and $\sigma$ computed on training set only, then applied to validation/test.

\section{S4. Validation Methodology}

\subsection{S4.1 Cross-Validation}

All modules use 5-fold stratified cross-validation to ensure robust performance estimates:

\[
\text{CV Score} = \frac{1}{5} \sum_{k=1}^{5} \text{Score}_k
\]

\subsection{S4.2 Uncertainty Quantification Calibration}

For Module 2 (Bayesian NN), calibration evaluated via:

\[
\text{Expected Calibration Error} = \frac{1}{M} \sum_{m=1}^{M} |\text{Accuracy}_m - \text{Confidence}_m|
\]

where $M$ is number of confidence bins. Target ECE $< 0.05$.

\section{S5. Computational Resources}

\begin{table}[H]
\centering
\caption{Training Computational Requirements}
\label{tab:compute}
\begin{tabular}{|l|l|l|l|}
\hline
\textbf{Module} & \textbf{Training Time} & \textbf{GPU} & \textbf{RAM} \\
\hline
Module 1 (GP) & 4-6 hours & CPU-only & 4 GB \\
\hline
Module 2 (BNN) & 1-2 hours & NVIDIA V100 & 8 GB \\
\hline
Module 3 (GNN) & 2-3 hours & NVIDIA V100 & 16 GB \\
\hline
Module 5 (CNN) & 8-12 hours & NVIDIA A100 & 24 GB \\
\hline
\textbf{Total} & \textbf{15-23 hours} & - & - \\
\hline
\end{tabular}
\end{table}

\section{S6. Inference Latency}

\begin{table}[H]
\centering
\caption{Deployment Inference Latency (ms per sample)}
\label{tab:latency}
\begin{tabular}{|l|l|l|l|}
\hline
\textbf{Module} & \textbf{CPU (ms)} & \textbf{GPU (ms)} & \textbf{Mobile (ms)} \\
\hline
Module 1 & 25-50 & 15-20 & N/A \\
\hline
Module 2 & 5-10 & 2-3 & 15-25 \\
\hline
Module 3 & 10-20 & 5-8 & 30-50 \\
\hline
Module 4 & 50-100 & 20-30 & 80-150 \\
\hline
Module 5 & 100-200 & 50-100 & 500-1000 \\
\hline
\end{tabular}
\end{table}

\section{S7. Error Analysis}

\subsection{S7.1 Module 1 Residual Analysis}

Residuals (predicted - actual) for phytochemical predictions follow approximately normal distribution:
- Mean residual: $-0.012\%$ (unbiased)
- Std dev: $0.503\%$
- Shapiro-Wilk test: $p=0.087$ (normal distribution not rejected)

\subsection{S7.2 Module 2 Bioavailability Prediction Errors}

Top sources of prediction error:
\begin{itemize}
    \item Gut microbiota composition (35\% variance explained) — not measured
    \item Concurrent medication interactions (25\%) — partially captured
    \item Individual metabolic rate variation (20\%)
    \item Food-herb interaction (15\%)
    \item Measurement noise in bioavailability studies (5\%)
\end{itemize}

\section{S8. Sensitivity Analysis}

\subsection{S8.1 Feature Importance (Module 2)}

Relative importance of patient features for bioavailability prediction (permutation importance):

\begin{table}[H]
\centering
\caption{Feature Importance for Bioavailability Prediction}
\label{tab:importance}
\begin{tabular}{|l|l|}
\hline
\textbf{Feature} & \textbf{Importance Score} \\
\hline
Delivery method (extract/powder/decoction/infusion) & 0.342 \\
\hline
Compound lipophilicity (LogP) & 0.187 \\
\hline
Patient age & 0.156 \\
\hline
Patient BMI & 0.134 \\
\hline
Gut health score & 0.108 \\
\hline
Molecular weight & 0.073 \\
\hline
\end{tabular}
\end{table}

\section{S9. Ablation Studies}

\subsection{S9.1 Module 1: Impact of ARIMA Component}

\begin{table}[H]
\centering
\caption{Module 1 Performance with and without ARIMA}
\label{tab:ablation_m1}
\begin{tabular}{|l|l|l|l|}
\hline
\textbf{Configuration} & \textbf{R²} & \textbf{RMSE (\%)} & \textbf{MAE (\%)} \\
\hline
GP only & 0.612 & 0.548 & 0.412 \\
\hline
GP + ARIMA & 0.639 & 0.504 & 0.382 \\
\hline
Improvement & +2.7\% & -7.6\% & -7.2\% \\
\hline
\end{tabular}
\end{table}

\subsection{S9.2 Module 3: Impact of Structural Features}

\begin{table}[H]
\centering
\caption{Module 3 Performance with and without Structural Features}
\label{tab:ablation_m3}
\begin{tabular}{|l|l|l|l|}
\hline
\textbf{Configuration} & \textbf{AUC} & \textbf{Sensitivity} & \textbf{Specificity} \\
\hline
DeepWalk embeddings only & 0.962 & 0.931 & 0.954 \\
\hline
Embeddings + struct features & 0.993 & 0.972 & 0.982 \\
\hline
Improvement & +3.1\% & +4.4\% & +2.9\% \\
\hline
\end{tabular}
\end{table}

\section{S10. Competitive Benchmarks}

\subsection{S10.1 Baseline Model Comparisons}

\begin{table}[H]
\centering
\caption{Module 1: Gaussian Process vs Baseline Regressors}
\label{tab:baseline_comp}
\small
\begin{tabular}{|l|l|l|l|}
\hline
\textbf{Model} & \textbf{R²} & \textbf{RMSE (\%)} & \textbf{Training Time (s)} \\
\hline
Linear Regression & 0.321 & 1.203 & 0.5 \\
\hline
Random Forest & 0.412 & 0.987 & 12.3 \\
\hline
XGBoost & 0.518 & 0.742 & 24.1 \\
\hline
Neural Network & 0.591 & 0.589 & 180 \\
\hline
Gaussian Process (ours) & 0.639 & 0.504 & 240 \\
\hline
\end{tabular}
\end{table}

\section{References}

\end{document}
